% Documento de herramientas de ingeniería (máximo 2 páginas).
\documentclass[10pt,a4paper]{article}
\usepackage[spanish,es-noshorthands,es-tabla]{babel}
\usepackage[utf8]{inputenc}
\usepackage[T1]{fontenc}
\usepackage{lmodern}
\usepackage[margin=1.6cm]{geometry}
\usepackage{booktabs,tabularx,xcolor,enumitem}
\usepackage{url}
\usepackage[numbers]{natbib}
\setlength{\bibsep}{0pt}
\setlist{nosep,leftmargin=*}
\renewcommand{\arraystretch}{1.1}

\title{\vspace{-1.5cm}Herramientas de ingeniería\\\large Proyecto Grupal 1: sistema distribuido de reducción adaptativa de ruido}
\author{Nicolás Florez Jiménez, Tamara Vanessa Cajiao Molina,\\Santiago Robles Obando, Santiago Boza Quiros\\CE 1110 Análisis de Señales Mixtas, II Semestre 2026}
\date{\vspace{-8mm}}

\begin{document}
\maketitle

\section{Selección de técnicas, recursos y herramientas}
Las herramientas se eligieron según las variables del problema: frecuencia de muestreo y
aliasing, latencia por bloque, memoria del microcontrolador, integridad del enlace y
posibilidad de verificar el firmware contra una referencia.

\noindent\small
\begin{tabularx}{\linewidth}{@{}>{\raggedright\arraybackslash}p{3.2cm}>{\raggedright\arraybackslash}p{3.6cm}>{\raggedright\arraybackslash}X@{}}
\toprule
Categoría & Herramienta o recurso & Variable que atiende \\
\midrule
Modelado y simulación & Python 3, NumPy, SciPy, Matplotlib~\cite{numpy,scipy}, Jupyter & Diseño FIR/IIR, DFT directa frente a FFT, métricas y figuras reproducibles con semillas fijas. \\
Firmware de adquisición & Pico SDK en C con ADC, DMA y TinyUSB~\cite{picosdk,rp2040} & Muestreo continuo a 16~kHz sin pérdidas, quitar DC, Hann y FFT de diagnóstico. \\
Firmware de filtrado & ESP-IDF 5.5 (C, FreeRTOS) desde PowerShell, SPI esclavo, DAC continuo, ESP Timer~\cite{espidf,esp32} & Dos rutas independientes en un microcontrolador de doble núcleo; tiempo por bloque medido en~$\mu$s. \\
Comunicación & SPI 2~MHz con trama propia y CRC-8~\cite{peterson}; USB serie a 1~Mbaud & Delimitación de bloques, detección de errores y pérdidas. \\
Datasheets & RP2040 (ADC), ESP32 (DAC 8 bits), MAX4466~\cite{max4466}, LM358~\cite{lm358}, nota de Sallen-Key~\cite{sloa024} & Rango de entrada 0--3.3~V, resolución, alimentación y diseño del filtro de reconstrucción. \\
Instrumentos & Multímetro digital ANENG DM850 & Alimentación de 5~V del LM358 y tensiones DC en sus pines durante la ejecución, de las que depende el filtro analógico. \\
Verificación & pytest, GCC, scripts de prueba en Python & Más de 150 pruebas en Python (modelo, procesador C frente al modelo Python, trama SPI, protocolo PC y cierre del DAC) y pruebas en C del receptor (DC, Hann, FFT frente a DFT, CRC-8). \\
Medición con la Pico & \emph{pico\_instrumento.py}, \emph{capturas\_pico.py}, \emph{demo analogico} & ADC1 por USB: 12 bloques separados por punto, amplitudes de GPIO25 frente a pin 7 del LM358. \\
Validación independiente & \emph{validar}, \emph{concordancia.py}; esptool & Mismo PCM16 contra Python, criterio digital $\geq90$\,\%; diagnóstico del puente USB. \\
Documentación y gestión & Git (master/development, ramas por función), Mermaid, \LaTeX{} IEEEtran & Trabajo incremental, diagramas versionados y artículo. \\
Inteligencia artificial & OpenAI Codex; Anthropic Claude & Desarrollo y depuración de bloques de código; investigación de fuentes y formato de documentación. \\
\bottomrule
\end{tabularx}
\normalsize

\section{Aplicación de las herramientas al problema}
El trabajo siguió una cadena en la que cada etapa se valida con la anterior:
\begin{enumerate}
\item \textbf{Modelo de referencia.} El cuaderno y \texttt{prototipo.py} generan los
casos con semillas, implementan la DFT por definición y la comparan con la FFT, diseñan
un FIR Kaiser y tres IIR, y calculan MSE, SNR, energía y tiempo. Los parámetros se leen
de archivos JSON, de modo que cada figura y tabla del artículo se regenera con un comando.
\item \textbf{Paso de Python a C.} \texttt{exportar\_coeficientes.py} conserva los 75
coeficientes FIR en float32. Las señales se transfieren desde CSV como PCM16 al
procesador de registros; las pruebas usan la misma entrada cuantizada que la placa.
La instrumentación agrega costos por ruta y latencia digital hasta aceptación UART;
la atenuación regional usa residuos y ganancias reales de la ESP32. La terminal
guarda notas y prepara audios de estímulo sin recortes para parlante--micrófono.
\item \textbf{Verificación cruzada.} Las pruebas compilan con GCC el mismo núcleo C que
corre en la ESP32 y lo comparan con Python: error por muestra $\leq3\times10^{-6}$ y
diferencias de métricas $\leq0.03$~dB. En la placa, la SNR de salida coincidió con
Python sobre la misma entrada PCM16: diferencia FIR $3.3\times10^{-5}$~dB y espectral
$8.2\times10^{-8}$~dB; concordancias de 100 y 99.99966\,\%, respectivamente.
\item \textbf{Cadena física y adquisición de datos.} Los datasheets fijaron el rango
del ADC (0--3.3~V, 12 bits) y del DAC (8 bits, GPIO25). Con el multímetro ANENG DM850 se verificó,
con el sistema en ejecución, que el LM358 recibe 5~V, condición para que el filtro
analógico funcione; además, la entrada no inversora (pin~3) y la salida (pin~1) del
primer amplificador marcaron 0.815~V. Para el filtro analógico, la Pico se reutilizó
como instrumento: \emph{pico\_instrumento.py} coordina USB y
\emph{capturas\_pico.py} interpreta ADC1. La demo mide los mismos tonos antes
(GPIO25) y después (pin~7) sin concatenar huecos. El modelo usa 22~nF y retención DAC;
los tres tonos útiles cumplen y 4~kHz queda inválido por variabilidad de 19.3\,\%.
El comando \emph{validar} reprocesa en la PC la misma
entrada guardada con los filtros de Python y calcula la concordancia con la salida de
la ESP32, sin reemplazar el procesamiento de la placa. Con esptool se comprobó que el
puente USB no respondía a una tasa superior, por lo que el enlace quedó en 1~Mbaud. La compilación, la programación y la ejecución se hicieron desde
PowerShell; la PC solo recibe y muestra salidas. Los scripts llevan los datos del circuito a
la PC: \texttt{graficar\_microfono.py} muestra en vivo la forma de onda y el espectro
que captura la Pico; \texttt{sesiones.py} configura parámetros sin recompilar, ordena
capturas y evaluaciones en la ESP32, guarda entrada, salida FIR y salida espectral como
WAV con \texttt{sesion.json}, métricas y registro de eventos, y verifica integridad con
CRC e índices. \emph{respuesta\_fir\_medida.py} compara los WAV de placa con los
coeficientes; \emph{resumen\_pruebas.py} produce medias, desviaciones y tabla \LaTeX{}
por condición a partir de las sesiones completas, como exige el informe de
repeticiones. No califica automáticamente el sistema. Registros completos locales
y tablas/figuras seleccionadas para la entrega se mantienen separados.
\item \textbf{Uso de IA.} Codex se usó para escribir y depurar bloques de código
(controladores, protocolo y pruebas), y Claude para buscar fuentes, revisar la
coherencia entre código y documentación y dar formato a los documentos. Toda salida se
aceptó solo después de revisarla, compilarla y pasar las pruebas automáticas; los
valores numéricos del artículo provienen de las ejecuciones, no de la IA.
\end{enumerate}

\section{Adaptación de las técnicas al problema}
\begin{itemize}
\item \textbf{FFT propia.} Se implementó una FFT radix-2 iterativa en C estándar, sin
dependencias del SDK, para poder probarla en la PC; se validó contra NumPy.
\item \textbf{FIR por bloques con estado.} El FIR se adaptó a un anillo de 75 muestras
que conserva el historial entre bloques de 256, y la comparación compensa su retardo
fijo de 37 muestras en lugar de buscar el desfase que mejor puntúe.
\item \textbf{Espectral con memoria acotada.} La superposición y suma usa solo tres
buffers de 1024 muestras; el umbral relativo se recalcula por bloque y protege la banda
útil. Sus parámetros cambian en ejecución, como exige la defensa.
\item \textbf{Protocolo propio.} La trama SPI (cabecera, identificador, 256 float32,
CRC-8) y la línea READY se diseñaron para que la ESP32 detecte errores, saltos y colas
llenas y aborte con evidencia, en vez de procesar datos corruptos.
\item \textbf{Herramientas de compilación.} Se creó \texttt{ejecutar\_idf.ps1}, que
compila desde una copia sin tildes porque ESP-IDF falla con rutas como
``generación\_señales''.
\item \textbf{Filtro analógico.} El Sallen-Key se dimensionó con la nota de TI para un
corte de 2283~Hz con 22~nF, el valor comercial disponible más cercano a los 20~nF de un Butterworth exacto.
\end{itemize}

{\small
\begin{thebibliography}{99}
\bibitem{numpy} C. R. Harris \emph{et al.}, ``Array programming with NumPy,'' \emph{Nature}, vol. 585, pp. 357--362, 2020.
\bibitem{scipy} P. Virtanen \emph{et al.}, ``SciPy 1.0: Fundamental algorithms for scientific computing in Python,'' \emph{Nature Methods}, vol. 17, pp. 261--272, 2020.
\bibitem{picosdk} Raspberry Pi Ltd., ``Raspberry Pi Pico-series C/C++ SDK.'' [En línea]. Disponible: \url{https://datasheets.raspberrypi.com/pico/raspberry-pi-pico-c-sdk.pdf}
\bibitem{rp2040} Raspberry Pi Ltd., ``RP2040 Datasheet,'' sec. 4.9. [En línea]. Disponible: \url{https://datasheets.raspberrypi.com/rp2040/rp2040-datasheet.pdf}
\bibitem{espidf} Espressif Systems, ``ESP-IDF Programming Guide v5.5.'' [En línea]. Disponible: \url{https://docs.espressif.com/projects/esp-idf/en/v5.5/esp32/}
\bibitem{esp32} Espressif Systems, ``ESP32 Series Datasheet,'' 2025.
\bibitem{peterson} W. W. Peterson y D. T. Brown, ``Cyclic codes for error detection,'' \emph{Proc. IRE}, vol. 49, no. 1, pp. 228--235, 1961.
\bibitem{max4466} Maxim Integrated, ``MAX4465--MAX4469 Micropower, Low-Noise Microphone Preamplifiers,'' datasheet, 2017.
\bibitem{lm358} Texas Instruments, ``LM358, LM258, LM2904: Industry-Standard Dual Operational Amplifiers,'' SLOS068AB, oct. 2024.
\bibitem{sloa024} J. Karki, ``Analysis of the Sallen-Key Architecture,'' Texas Instruments, SLOA024B, 2002.
\end{thebibliography}}

\end{document}
